% !TEX program = pdflatex
\documentclass[11pt]{article}

% ==========================
% Encoding / language (pdflatex)
% ==========================
\usepackage[utf8]{inputenc}
\usepackage[T1]{fontenc}
\usepackage[english]{babel}
\usepackage{lmodern}
\usepackage{microtype}

% In the preamble
\usepackage{amsthm}

% Numbered within sections (recommended)
\newtheorem{theorem}{Theorem}[section]
\newtheorem{proposition}[theorem]{Proposition}
\newtheorem{lemma}[theorem]{Lemma}
\newtheorem{corollary}[theorem]{Corollary}

\theoremstyle{definition}
\newtheorem{definition}[theorem]{Definition}
\newtheorem{exercise}[theorem]{Exercise}

\theoremstyle{remark}
\newtheorem{remark}[theorem]{Remark}
% ==========================
% Page layout
% ==========================
\usepackage[margin=1in]{geometry}

% ==========================
% Math
% ==========================
\usepackage{amsmath,amssymb,amsthm,mathtools}
\usepackage{enumitem}

% ==========================
% Macros
% ==========================
\newcommand{\Pbb}{\mathbb{P}}
\newcommand{\Qbb}{\mathbb{Q}}
\newcommand{\E}{\mathbb{E}}
\newcommand{\R}{\mathbb{R}}
\newcommand{\Q}{\mathbb{Q}}
\newcommand{\F}{\mathcal{F}}
\newcommand{\1}{\mathbf{1}}
\newcommand{\dd}{\,\mathrm{d}}
\newcommand{\wt}{\widetilde}
\newcommand{\bs}{\boldsymbol}

% ==========================
% Theorem-like
% ==========================
\theoremstyle{definition}

% ==========================
% "Board" environment (English)
% ==========================
\newenvironment{board}{%
  \par\medskip
  \noindent
  \fbox{%
    \begin{minipage}{0.98\linewidth}
    \textbf{On the board (plan + comments).}\par\medskip
}{%
    \end{minipage}%
  }%
  \par\medskip
}

% ==========================
% Document
% ==========================
\begin{document}

\begin{center}
{\Large \textbf{Lecture notes --- It\^o examples \& first steps toward Girsanov}}\\[0.25em]
{\small (Board-ready derivations; focus: explicit computations + intuition for change of measure)}
\end{center}

\vspace{-0.5em}


% ==========================================================
\section*{Part A --- It\^o formula: 40--45 minutes of explicit computations}
% ==========================================================

\subsection*{A.1 Warm-up: the prototype computation}
Let $(W_t)_{t\ge 0}$ be a Brownian motion and recall It\^o's formula:
for $f\in C^{1,2}([0,T]\times\R)$ and an It\^o process $X_t$ with
\[
\dd X_t = b_t\,\dd t + \sigma_t\,\dd W_t,
\]
we have
\[
\dd f(t,X_t)
=
\partial_t f(t,X_t)\,\dd t
+ \partial_x f(t,X_t)\,\dd X_t
+ \frac12 \partial_{xx} f(t,X_t)\,(\dd X_t)^2,
\qquad (\dd X_t)^2=\sigma_t^2\,\dd t.
\]


\begin{itemize}[leftmargin=1.5em, itemsep=0.25em]
\item Always separate: \textbf{(drift)} $\dd t$ part vs \textbf{(martingale)} $\dd W$ part.
\item Remember: $(\dd W_t)^2=\dd t$, $\dd t\,\dd W_t=0$, $(\dd t)^2=0$.
\end{itemize}


\subsection*{A.2 Example 1: $W_t^2-t$ is a martingale}
Apply It\^o to $f(x)=x^2$ with $X_t=W_t$.
We have $f'(x)=2x$, $f''(x)=2$, and $\dd W_t = 0\cdot \dd t + 1\cdot \dd W_t$:
\[
\dd(W_t^2) = 2W_t\,\dd W_t + \frac12\cdot 2 \,(\dd W_t)^2
= 2W_t\,\dd W_t + \dd t.
\]
Hence
\[
\dd(W_t^2-t)=2W_t\,\dd W_t,
\]
so $W_t^2-t$ is an It\^o integral and therefore a martingale (integrable, mean $0$ increment).


\textbf{What to highlight.} The \emph{only} reason we need the ``$-t$'' is the It\^o correction term.


\subsection*{A.3 Example 2: Exponential martingale (``kill the drift'')}
Let $Y_t=\exp(a t + b W_t)$, with constants $a,b\in\R$.
Use It\^o with $f(t,x)=e^{at+bx}$.
Compute derivatives:
\[
\partial_t f = a f,\qquad \partial_x f = b f,\qquad \partial_{xx}f=b^2 f.
\]
With $X_t=W_t$:
\[
\dd Y_t = aY_t\,\dd t + bY_t\,\dd W_t + \tfrac12 b^2 Y_t\,\dd t
= Y_t\Big( (a+\tfrac12 b^2)\,\dd t + b\,\dd W_t\Big).
\]
Choose $a=-\tfrac12 b^2$ to remove the drift:
\[
\dd Y_t = bY_t\,\dd W_t,
\quad\Rightarrow\quad
Y_t=\exp\Big(bW_t-\tfrac12 b^2 t\Big)
\text{ is a martingale.}
\]


\textbf{Pedagogical link.} This is the prototype of the Radon--Nikodym density used in Girsanov.


\subsection*{A.4 Example 3: Geometric Brownian motion and the log trick}
Consider GBM:
\[
\dd S_t = \mu S_t\,\dd t + \sigma S_t\,\dd W_t,\qquad S_0>0.
\]
Let $f(x)=\ln x$, $f'(x)=1/x$, $f''(x)=-1/x^2$.
Then
\[
\dd(\ln S_t)
=
\frac{1}{S_t}\dd S_t + \frac12\Big(-\frac{1}{S_t^2}\Big)(\dd S_t)^2
=
\frac{1}{S_t}(\mu S_t\,\dd t+\sigma S_t\,\dd W_t)
-\frac12\frac{1}{S_t^2}(\sigma^2 S_t^2\,\dd t),
\]
so
\[
\dd(\ln S_t) = \Big(\mu-\tfrac12\sigma^2\Big)\dd t + \sigma\,\dd W_t.
\]
Integrate:
\[
\ln S_t = \ln S_0 + \Big(\mu-\tfrac12\sigma^2\Big)t + \sigma W_t,
\qquad
S_t = S_0\exp\Big(\Big(\mu-\tfrac12\sigma^2\Big)t+\sigma W_t\Big).
\]


\textbf{Board payoff.} You get the explicit solution + lognormal law in one shot.


\subsection*{A.5 Example 4: Discounting trick (product rule, no PDE yet)}
Let $A_t$ be a deterministic (or finite-variation adapted) process with
\[
\dd A_t = -\gamma_t A_t\,\dd t,
\qquad A_0=1,
\]
and let $S_t$ satisfy $\dd S_t = \mu_t S_t\dd t + \sigma_t S_t\dd W_t$.
Define $X_t:=A_t S_t$. Use product rule:
\[
\dd X_t = A_t\,\dd S_t + S_t\,\dd A_t + \dd[A,S]_t.
\]
Here $A$ has zero quadratic variation, so $\dd[A,S]_t=0$, hence
\[
\dd X_t = A_t(\mu_t S_t\,\dd t+\sigma_t S_t\,\dd W_t) - \gamma_t A_t S_t\,\dd t
= X_t(\mu_t-\gamma_t)\dd t + X_t\sigma_t\,\dd W_t.
\]
Choose $\gamma_t=\mu_t$:
\[
\dd X_t = X_t\sigma_t\,\dd W_t,
\]
so $(X_t)$ is a martingale (under suitable integrability, e.g. $\E\int_0^T X_t^2\sigma_t^2\dd t<\infty$).


\textbf{Message.} ``Find the right discount factor $\Rightarrow$ spot the right drift for pricing.''


\subsection*{A.6 (If time) Example 5: A simple It\^o identity to remove a stochastic integral}
Goal: compute an expression for $\int_0^t e^{W_s}\,\dd W_s$ \emph{without} stochastic integral.

Take $f(x)=e^x$ and $X_t=W_t$:
\[
\dd(e^{W_t}) = e^{W_t}\,\dd W_t + \tfrac12 e^{W_t}\,\dd t.
\]
Integrate from $0$ to $t$:
\[
e^{W_t}-1 = \int_0^t e^{W_s}\,\dd W_s + \frac12\int_0^t e^{W_s}\,\dd s,
\]
so
\[
\boxed{
\int_0^t e^{W_s}\,\dd W_s
=
e^{W_t}-1-\frac12\int_0^t e^{W_s}\,\dd s.}
\]

\subsection*{A.7: Norm squared $\|W_t\|^2$ and drift identification}
Assume for simplicity $\Sigma=I_d$ (independent coordinates).
Let $R_t:=\|W_t\|^2=\sum_{i=1}^d (W_t^i)^2$.
Using $\dd (W^i)^2 = 2W^i\,\dd W^i + \dd t$ and summing:
\[
\dd R_t = 2\sum_{i=1}^d W_t^i\,\dd W_t^i + d\,\dd t
= 2\,W_t^\top \dd W_t + d\,\dd t.
\]
\textbf{Extract:} $\|W_t\|^2-dt$ is a martingale.
\textbf{Comment:} this shows the ``dimension drift'' $+d\,dt$.

% ----------------------------------------------------------
\subsection*{A.8: MD-8: It\^o for $f(X^1_t,X^2_t)$ with correlated Brownian drivers}


\textbf{Setup (correlated Brownian motion).}
Let $(W_t^1,W_t^2)$ be a 2D Brownian motion with correlation $\rho\in[-1,1]$:
\[
\dd[W^1]_t=\dd t,\qquad \dd[W^2]_t=\dd t,\qquad \dd[W^1,W^2]_t=\rho\,\dd t.
\]
Let $X^1,X^2$ be It\^o processes driven by $(W^1,W^2)$:
\[
\dd X_t^1 = b_1(t)\,\dd t + \sigma_{11}(t)\,\dd W_t^1 + \sigma_{12}(t)\,\dd W_t^2,
\qquad
\dd X_t^2 = b_2(t)\,\dd t + \sigma_{21}(t)\,\dd W_t^1 + \sigma_{22}(t)\,\dd W_t^2.
\]
Let $f\in C^2(\R^2)$ and set $Y_t:=f(X_t^1,X_t^2)$.

\medskip
\textbf{Step 1: write the quadratic covariations.}
Using bilinearity and $\dd W^1\dd W^2=\rho\,\dd t$:
\[
\dd[X^1]_t
=
\Big(\sigma_{11}^2 + \sigma_{12}^2 + 2\rho\,\sigma_{11}\sigma_{12}\Big)\dd t,
\]
\[
\dd[X^2]_t
=
\Big(\sigma_{21}^2 + \sigma_{22}^2 + 2\rho\,\sigma_{21}\sigma_{22}\Big)\dd t,
\]
\[
\dd[X^1,X^2]_t
=
\Big(\sigma_{11}\sigma_{21}+\sigma_{12}\sigma_{22}
+\rho(\sigma_{11}\sigma_{22}+\sigma_{12}\sigma_{21})\Big)\dd t.
\]

\medskip
\textbf{Step 2: apply 2D It\^o formula.}
\[
\dd Y_t
=
f_{x_1}(X_t)\,\dd X_t^1
+
f_{x_2}(X_t)\,\dd X_t^2
+\frac12 f_{x_1x_1}(X_t)\,\dd[X^1]_t
+ f_{x_1x_2}(X_t)\,\dd[X^1,X^2]_t
+\frac12 f_{x_2x_2}(X_t)\,\dd[X^2]_t.
\]

\textbf{Board takeaway:} the correlation $\rho$ appears only through $\dd[X^1]$, $\dd[X^2]$, and especially $\dd[X^1,X^2]$,
which multiplies the mixed derivative $f_{x_1x_2}$.

% ----------------------------------------------------------
\subsection*{Concrete choice: $f(x_1,x_2)=x_1x_2$ (product)}
Take $f(x_1,x_2)=x_1x_2$. Then
\[
f_{x_1}=x_2,\quad f_{x_2}=x_1,\quad f_{x_1x_2}=1,\quad f_{x_1x_1}=f_{x_2x_2}=0.
\]
Hence
\[
\dd(X_t^1X_t^2)
=
X_t^2\,\dd X_t^1 + X_t^1\,\dd X_t^2 + \dd[X^1,X^2]_t.
\]
Plug the formula for $\dd[X^1,X^2]_t$ above to see explicitly where $\rho$ enters.

% ----------------------------------------------------------
\subsection*{Concrete choice: $f(x_1,x_2)=\exp(\alpha x_1+\beta x_2)$ (drift-killing style)}
Let $f(x_1,x_2)=\exp(\alpha x_1+\beta x_2)$, so
\[
f_{x_1}=\alpha f,\quad f_{x_2}=\beta f,\quad
f_{x_1x_1}=\alpha^2 f,\quad f_{x_2x_2}=\beta^2 f,\quad f_{x_1x_2}=\alpha\beta f.
\]
Then
\[
\frac{\dd Y_t}{Y_t}
=
\alpha\,\dd X_t^1 + \beta\,\dd X_t^2
+\frac12\alpha^2\,\dd[X^1]_t
+\alpha\beta\,\dd[X^1,X^2]_t
+\frac12\beta^2\,\dd[X^2]_t.
\]
\textbf{Board takeaway:} compared to 1D, the extra term is the cross-variation piece
$\alpha\beta\,\dd[X^1,X^2]_t$, which carries the correlation $\rho$.
% ==========================================================

% ==========================================================
\section*{Part B --- Black--Scholes hedging $\Rightarrow$ PDE $\Rightarrow$ expectation}
% ==========================================================

% ----------------------------------------------------------
\subsection*{B.1 Model setup (GBM + bank account)}
% ----------------------------------------------------------
Assume a frictionless market with one risky asset $S$ and one money-market account $B$:
\[
\frac{\dd S_t}{S_t}=\mu\,\dd t+\sigma\,\dd W_t,
\qquad \mu\in\R,\ \sigma>0,
\]
\[
\dd B_t=rB_t\,\dd t,\qquad B_0=1,\qquad B_t=e^{rt}.
\]
Define the discount factor
\[
D_t:=B_t^{-1}=e^{-rt}.
\]

\begin{itemize}[leftmargin=1.6em, itemsep=0.2em]
\item Parameters: $\mu$ (real drift), $\sigma$ (volatility), $r$ (short rate).
\item $B_t=e^{rt}$ is deterministic; $D_t=e^{-rt}$ is its inverse.
\end{itemize}

% ----------------------------------------------------------
\subsection*{B.2 Derivative price as a Markov function and It\^o expansion}
% ----------------------------------------------------------
Fix a maturity $T$ and a European payoff $g(S_T)$.
Assume the time-$t$ price is given by a smooth deterministic function
\[
V_t=p(t,S_t),
\qquad 0\le t\le T,
\qquad p(T,s)=g(s).
\]
By It\^o's formula (with $(\dd S_t)^2=\sigma^2 S_t^2\,\dd t$),
\[
\dd p(t,S_t)
=
\Big(\partial_t p(t,S_t)
+\mu S_t\,\partial_s p(t,S_t)
+\tfrac12\sigma^2 S_t^2\,\partial_{ss}p(t,S_t)\Big)\dd t
+\sigma S_t\,\partial_s p(t,S_t)\,\dd W_t.
\]

\begin{itemize}[leftmargin=1.6em, itemsep=0.2em]
\item Write explicitly: \emph{drift part} + \emph{noise part}.
\item The coefficient of $\dd W_t$ is $\sigma S_t\,\partial_s p(t,S_t)$.
\end{itemize}

% ----------------------------------------------------------
\subsection*{B.3 Self-financing hedging portfolio and ``kill the $\dd W$''}
% ----------------------------------------------------------
Consider a self-financing portfolio holding $\Delta_t$ shares of stock and the rest in the bank account:
\[
X_t := \Delta_t S_t + \beta_t B_t.
\]
Self-financing means ``no external cash injection'':
\[
\dd X_t = \Delta_t\,\dd S_t + \beta_t\,\dd B_t.
\]
Using $\beta_t B_t = X_t-\Delta_t S_t$ and $\dd B_t=rB_t\,\dd t$,
\[
\dd X_t
=
\Delta_t(\mu S_t\,\dd t+\sigma S_t\,\dd W_t)
+r\big(X_t-\Delta_t S_t\big)\dd t.
\]

We aim to replicate the option: $X_t \equiv p(t,S_t)$, so $\dd X_t=\dd p(t,S_t)$.
Comparing the \textbf{$\dd W_t$ coefficients} gives the \textbf{delta hedge}:
\[
\boxed{\ \Delta_t=\partial_s p(t,S_t)\ }.
\]

\begin{itemize}[leftmargin=1.6em, itemsep=0.2em]
\item This is the ``kill the noise'' step: choose $\Delta_t$ to cancel the $\dd W_t$ term.
\item Once the noise is removed, the hedged position becomes locally riskless.
\end{itemize}

% ----------------------------------------------------------
\subsection*{B.4 Black--Scholes PDE (no-arbitrage)}
% ----------------------------------------------------------
Plug $\Delta_t=\partial_s p(t,S_t)$ into the drift matching.
From the portfolio drift:
\[
\dd X_t
=
\Big(rX_t+(\mu-r)S_t\Delta_t\Big)\dd t
+\sigma S_t\Delta_t\,\dd W_t.
\]
From the It\^o expansion:
\[
\dd p(t,S_t)
=
\Big(\partial_t p+\mu S_t\partial_s p+\tfrac12\sigma^2 S_t^2\partial_{ss}p\Big)\dd t
+\sigma S_t\partial_s p\,\dd W_t.
\]
After setting $\Delta_t=\partial_s p(t,S_t)$, the $\dd W_t$ terms match automatically.
Matching the drift terms yields
\[
r\,p(t,S_t)+(\mu-r)S_t\partial_s p(t,S_t)
=
\partial_t p(t,S_t)+\mu S_t\partial_s p(t,S_t)
+\tfrac12\sigma^2 S_t^2\partial_{ss}p(t,S_t).
\]
The $\mu S_t\partial_s p$ cancels on both sides, giving the \textbf{valuation PDE}:
\[
\boxed{\ 
\partial_t p(t,s)+rs\,\partial_s p(t,s)
+\tfrac12\sigma^2 s^2\,\partial_{ss}p(t,s)-r\,p(t,s)=0,
\qquad p(T,s)=g(s).
\ }
\]

\begin{itemize}[leftmargin=1.6em, itemsep=0.2em]
\item \textbf{Key point: $\mu$ disappears.} Prices do not depend on the real drift.
\item This is the core Black--Scholes hedging argument.
\end{itemize}

% ==========================================================
\section*{Solving the Black--Scholes PDE via the heat equation}
% ==========================================================

We now solve
\[
\partial_t p + rs\,\partial_s p + \tfrac12\sigma^2 s^2\,\partial_{ss}p - rp = 0,
\qquad p(T,s)=g(s).
\]

% ----------------------------------------------------------
\subsection*{Step 1: time-to-maturity and discounting}
% ----------------------------------------------------------
Let
\[
\tau:=T-t,
\qquad v(\tau,s):=p(T-\tau,s),
\qquad v(0,s)=g(s).
\]
Since $\partial_t p = -\partial_\tau v$, the PDE becomes
\[
\partial_\tau v
=
rs\,\partial_s v + \tfrac12\sigma^2 s^2\,\partial_{ss}v - r v.
\]
Remove the $-rv$ term via discounting:
\[
u(\tau,s):=e^{r\tau}v(\tau,s).
\]
A short computation gives
\[
\partial_\tau u
=
rs\,\partial_s u + \tfrac12\sigma^2 s^2\,\partial_{ss}u,
\qquad u(0,s)=g(s).
\]

% ----------------------------------------------------------
\subsection*{Step 2: logarithmic change of variables}
% ----------------------------------------------------------
Set
\[
x:=\ln s,\qquad w(\tau,x):=u(\tau,e^x).
\]
Using
\[
\partial_s u=\frac{1}{s}w_x,\qquad
\partial_{ss}u=\frac{1}{s^2}(w_{xx}-w_x),
\]
we obtain
\[
\partial_\tau w
=
\tfrac12\sigma^2 w_{xx} + \Big(r-\tfrac12\sigma^2\Big)w_x,
\qquad w(0,x)=g(e^x).
\]

% ----------------------------------------------------------
\subsection*{Step 3: remove the transport term (shift)}
% ----------------------------------------------------------
Let
\[
a:=r-\tfrac12\sigma^2,
\qquad
\varphi(\tau,y):=w(\tau,y-a\tau).
\]
Then $\varphi$ solves the heat equation
\[
\partial_\tau \varphi
=
\tfrac12\sigma^2\,\partial_{yy}\varphi,
\qquad
\varphi(0,y)=g(e^y).
\]

% ----------------------------------------------------------
\subsection*{Step 4: heat equation solution $\Rightarrow$ expectation form}
% ----------------------------------------------------------
If $Z\sim\mathcal N(0,1)$, the heat equation solution is
\[
\varphi(\tau,y)
=
\E\!\left[g\!\left(e^{\,y+\sigma\sqrt{\tau}\,Z}\right)\right].
\]
Undoing the changes of variables yields, for $\tau=T-t$,
\[
\boxed{
\;
p(t,s)
=
e^{-r(T-t)}
\E\!\left[
g\!\left(
s\exp\!\Big((r-\tfrac12\sigma^2)(T-t)+\sigma\sqrt{T-t}\,Z\Big)
\right)
\right].
\;}
\]

\begin{remark}[Pedagogical takeaway]
The expectation representation comes from \emph{solving the PDE} (via the heat equation).
Next, Girsanov will explain how the same representation arises from a change of probability measure.
\end{remark}

% ==========================================================
\section*{Closed-form call formula and hedging}
% ==========================================================

From now on, take $g(s)=(s-K)_+$ (European call with strike $K$).
A direct computation of the expectation above yields the Black--Scholes formula.

\begin{theorem}[Black--Scholes call price]
\[
\boxed{
\;
p(0,S_0)=S_0\,\mathbf N(d_+)-K e^{-rT}\,\mathbf N(d_-),
\;}
\qquad
d_\pm
=
\frac{\ln\!\big(\frac{S_0}{K e^{-rT}}\big)}{\sigma\sqrt{T}}
\pm \frac12\sigma\sqrt{T}.
\]
\end{theorem}

\begin{theorem}[Delta hedge]
\[
\boxed{
\;
\Delta_t=\partial_s p(t,S_t)=\mathbf N\!\big(d_+(t,S_t)\big),
\;}
\qquad
d_+(t,s)
=
\frac{\ln\!\big(\frac{s}{K e^{-r(T-t)}}\big)}{\sigma\sqrt{T-t}}
+\frac12\sigma\sqrt{T-t}.
\]
\end{theorem}

% ----------------------------------------------------------
\subsection*{Greeks (definitions + interpretation)}
% ----------------------------------------------------------
Define the main sensitivities:
\[
\Delta := \partial_s p,\qquad
\Gamma := \partial_{ss}p,\qquad
\Theta := \partial_t p,\qquad
\text{Vega}:=\partial_\sigma p,\qquad
\rho:=\partial_r p.
\]
\begin{itemize}[leftmargin=1.6em, itemsep=0.2em]
\item $\Delta$: hedge ratio (shares in the replicating strategy).
\item $\Gamma$: curvature; drives rebalancing intensity (liquidity/transaction-cost sensitivity).
\item $\Theta$: time decay of the option value.
\item Vega: sensitivity to volatility misspecification (positive for calls).
\item $\rho$: sensitivity to the interest rate level.
\end{itemize}

\begin{remark}
In practice, ``model risk'' is precisely the sensitivity of prices to parameters:
the Greeks are first-order diagnostics of how wrong-model inputs distort valuation.
\end{remark}


% ==========================================================
% \section*{Part B --- Black--Scholes hedging $\Rightarrow$ PDE $\Rightarrow$ expectation}
% % ==========================================================

% \subsection*{B.1 Model setup (GBM + bank account)}
% Assume a frictionless market with:
% \[
% \frac{\dd S_t}{S_t}=\mu\,\dd t+\sigma\,\dd W_t,
% \qquad \mu\in\R,\ \sigma>0,
% \]
% and a money market account with constant short rate $r$:
% \[
% \dd B_t=rB_t\,\dd t,\qquad B_0=1,\qquad B_t=e^{rt}.
% \]
% Let $D_t:=B_t^{-1}=e^{-rt}$ denote the discount factor.


% \begin{itemize}
% \item Recall: GBM dynamics and interpretation of parameters $(\mu,\sigma,r)$.
% \item Write $B_t=e^{rt}$ and $D_t=e^{-rt}$.
% \end{itemize}


% \subsection*{B.2 Derivative price as a function $p(t,s)$ and It\^o expansion}
% Consider a European payoff $g(S_T)$ at maturity $T$.
% Assume its time-$t$ price is given by a deterministic function
% \[
% V_t=p(t,S_t),\qquad 0\le t\le T,\qquad p(T,s)=g(s).
% \]
% By It\^o's formula (1D),
% \[
% \dd p(t,S_t)
% =
% \partial_t p(t,S_t)\,\dd t
% +\partial_s p(t,S_t)\,\dd S_t
% +\frac12 \partial_{ss}p(t,S_t)\,(\sigma^2 S_t^2)\,\dd t.
% \]


% \begin{itemize}
% \item Write It\^o carefully: $(\dd S_t)^2=\sigma^2 S_t^2\,\dd t$.
% \item Separate drift part and $\dd W$ part explicitly.
% \end{itemize}


% \subsection*{B.3 Self-financing hedging portfolio and ``kill the $\dd W$''}
% The seller holds $\Delta_t$ shares of stock and invests the remaining wealth in the bank account.
% Define the portfolio value
% \[
% X_t := \Delta_t S_t + \beta_t B_t,
% \]
% and assume the strategy is self-financing:
% \[
% \dd X_t = \Delta_t\,\dd S_t + \beta_t\,\dd B_t
% = \Delta_t\,\dd S_t + r\big(X_t-\Delta_t S_t\big)\,\dd t.
% \]
% We aim to replicate the option: $X_t \equiv p(t,S_t)$, hence $\dd X_t=\dd p(t,S_t)$.
% Comparing the \textbf{$\dd W_t$ coefficients} in $\dd p(t,S_t)$ and $\dd X_t$,
% choose the \textbf{delta hedge}
% \[
% \boxed{\ \Delta_t=\partial_s p(t,S_t)\ }.
% \]


% \begin{itemize}
% \item Write $\dd X_t = \Delta_t(\mu S_t\,\dd t+\sigma S_t\,\dd W_t)+r(X_t-\Delta_t S_t)\dd t$.
% \item Match the $\dd W$ terms: $\sigma S_t \Delta_t = \sigma S_t \partial_s p$.
% \item Conclude: $\Delta_t=\partial_s p(t,S_t)$.
% \end{itemize}


% \subsection*{B.4 Derive the Black--Scholes PDE (no-arbitrage)}
% With $\Delta_t=\partial_s p(t,S_t)$, the stochastic term is hedged away and
% \[
% \dd\big(X_t-p(t,S_t)\big)=0
% \quad\Rightarrow\quad
% \dd X_t = \dd p(t,S_t).
% \]
% Using the drift parts only:
% \[
% r\big(p(t,S_t)-S_t\partial_s p(t,S_t)\big)\,\dd t
% =
% \partial_t p(t,S_t)\,\dd t
% +\mu S_t\partial_s p(t,S_t)\,\dd t
% +\frac12\sigma^2 S_t^2 \partial_{ss}p(t,S_t)\,\dd t
% -\mu S_t\partial_s p(t,S_t)\,\dd t,
% \]
% so the $\mu$ terms cancel (key point) and we obtain the \textbf{valuation PDE}:
% \[
% \boxed{\ \partial_t p(t,s)+rs\,\partial_s p(t,s)
% +\frac12\sigma^2 s^2\,\partial_{ss}p(t,s)-r\,p(t,s)=0,\qquad p(T,s)=g(s). \ }
% \]


% \begin{itemize}
% \item Emphasize: \textbf{$\mu$ disappears} $\Rightarrow$ price does not depend on real drift.
% \item This is exactly the Black--Scholes hedging argument in the slides. 
% \end{itemize}

% % ==========================================================
% \section{Solving the Black--Scholes PDE via the heat equation}
% % ==========================================================

% We consider the Black--Scholes valuation PDE (constant coefficients):
% \[
% \partial_t p(t,s)
% + r s\,\partial_s p(t,s)
% + \frac12\sigma^2 s^2\,\partial_{ss}p(t,s)
% - r\,p(t,s)=0,
% \qquad
% p(T,s)=g(s).
% \]

% \begin{remark}
% The key conceptual point (already emphasized) is that the real drift $\mu$ does not appear:
% the price depends only on $(r,\sigma)$.
% \end{remark}

% % ----------------------------------------------------------
% \subsection{Step 1: Time reversal and discounting}
% % ----------------------------------------------------------

% Introduce the time-to-maturity variable
% \[
% \tau := T-t,
% \qquad
% v(\tau,s):=p(T-\tau,s).
% \]
% Then $\partial_t p=-\partial_\tau v$, and the PDE becomes
% \[
% \partial_\tau v
% = r s\,\partial_s v
% + \frac12\sigma^2 s^2\,\partial_{ss}v
% - r v,
% \qquad
% v(0,s)=g(s).
% \]

% To eliminate the term $-rv$, define the discounted function
% \[
% u(\tau,s):=e^{r\tau}v(\tau,s).
% \]
% A direct computation yields
% \[
% \partial_\tau u
% = r s\,\partial_s u
% + \frac12\sigma^2 s^2\,\partial_{ss}u,
% \qquad
% u(0,s)=g(s).
% \]

% % ----------------------------------------------------------
% \subsection{Step 2: Logarithmic change of variables}
% % ----------------------------------------------------------

% Set
% \[
% x:=\ln s,
% \qquad
% w(\tau,x):=u(\tau,e^x).
% \]
% Using
% \[
% \partial_s u=\frac{1}{s}w_x,
% \qquad
% \partial_{ss}u=\frac{1}{s^2}(w_{xx}-w_x),
% \]
% the PDE becomes
% \[
% \partial_\tau w
% =
% \frac12\sigma^2 w_{xx}
% +\Big(r-\tfrac12\sigma^2\Big) w_x,
% \qquad
% w(0,x)=g(e^x).
% \]

% % ----------------------------------------------------------
% \subsection{Step 3: Removing the transport term}
% % ----------------------------------------------------------

% Define
% \[
% a:=r-\tfrac12\sigma^2,
% \qquad
% y:=x+a\tau,
% \qquad
% \varphi(\tau,y):=w(\tau,y-a\tau).
% \]
% Then $\varphi$ satisfies the heat equation
% \[
% \partial_\tau \varphi
% =
% \frac12\sigma^2\,\partial_{yy}\varphi,
% \qquad
% \varphi(0,y)=g(e^y).
% \]

% % ----------------------------------------------------------
% \subsection{Step 4: Solution of the heat equation}
% % ----------------------------------------------------------

% The solution is given by convolution with the Gaussian kernel:
% \[
% \varphi(\tau,y)
% =
% \int_{\mathbb{R}}
% g(e^z)\,
% \frac{1}{\sqrt{2\pi\sigma^2\tau}}
% \exp\!\left(-\frac{(y-z)^2}{2\sigma^2\tau}\right)\,\dd z.
% \]

% Equivalently, if $Z\sim\mathcal N(0,1)$,
% \[
% \varphi(\tau,y)
% =
% \E\!\left[g\!\left(e^{\,y+\sigma\sqrt{\tau}\,Z}\right)\right].
% \]

% % ----------------------------------------------------------
% \subsection{Step 5: Back to the price function}
% % ----------------------------------------------------------

% Undoing all changes of variables, with $\tau=T-t$,
% \[
% p(t,s)
% =
% e^{-r(T-t)}
% \,
% \E\!\left[
% g\!\left(
% s\exp\!\left(
% (r-\tfrac12\sigma^2)(T-t)
% +\sigma\sqrt{T-t}\,Z
% \right)
% \right)
% \right].
% \]

% This yields the pricing representation
% \[
% \boxed{
% \;
% p(t,s)
% =
% \E^{\Q}\!\left[
% e^{-r(T-t)}\,g(S_T)
% \mid S_t=s
% \right],
% \;}
% \]
% where under $\Q$,
% \[
% S_T
% =
% s\exp\!\left(
% (r-\tfrac12\sigma^2)(T-t)
% +\sigma(W_T-W_t)
% \right).
% \]

% \begin{remark}
% This representation will later be justified rigorously by
% martingale pricing and Girsanov's theorem.
% \end{remark}

% % ==========================================================
% \section{Closed-form Black--Scholes formula and hedging}
% % ==========================================================

% We now specialize to the classical Black--Scholes setting with
% \emph{constant} parameters:
% \[
% \dd S_t = \mu S_t\,\dd t + \sigma S_t\,\dd W_t,
% \qquad
% r,\sigma>0 \ \text{constants}.
% \]

% We consider a European call option with maturity $T$ and strike $K$:
% \[
% H=(S_T-K)_+.
% \]

% % ----------------------------------------------------------
% \subsection{Risk-neutral pricing formula}
% % ----------------------------------------------------------

% From the previous section, we know that the option price can be written as
% \[
% p(0,S_0)
% =
% \E^{\Q}\!\left[
% e^{-rT}(S_T-K)_+
% \right],
% \]
% where under the risk-neutral measure $\Q$,
% \[
% S_T
% =
% S_0
% \exp\!\left(
% \Big(r-\tfrac12\sigma^2\Big)T
% +\sigma\sqrt{T}\,Z
% \right),
% \qquad Z\sim\mathcal N(0,1).
% \]

% A direct computation of this expectation yields the celebrated
% \textbf{Black--Scholes closed-form formula}.

% \begin{theorem}[Black--Scholes call price]
% The time-$0$ price of a European call option is
% \[
% \boxed{
% \;
% p(0,S_0)
% =
% S_0\,\mathbf N(d_+)
% -
% K e^{-rT}\,\mathbf N(d_-),
% \;}
% \]
% where
% \[
% \mathbf N(x):=\int_{-\infty}^x \frac{e^{-y^2/2}}{\sqrt{2\pi}}\,\dd y
% \]
% is the standard normal distribution function, and
% \[
% d_\pm
% =
% \frac{\ln\!\left(\frac{S_0}{K e^{-rT}}\right)}{\sigma\sqrt{T}}
% \pm
% \frac12\sigma\sqrt{T}.
% \]
% \end{theorem}

% \begin{remark}[Interpretation of the terms]
% \begin{itemize}[leftmargin=2em]
% \item $S_0$ is the current stock price.
% \item $K e^{-rT}$ is the \emph{discounted strike}.
% \item $\sigma^2 T$ is the total variance over $[0,T]$.
% \item $\mathbf N(d_+)$ and $\mathbf N(d_-)$ are probabilities under $\Q$
% associated with being in-the-money.
% \end{itemize}
% \end{remark}

% % ----------------------------------------------------------
% \subsection{Hedging strategy (Delta)}
% % ----------------------------------------------------------

% The optimal self-financing hedging strategy is obtained by differentiation
% of the price with respect to the stock price.

% \begin{theorem}[Delta hedge]
% The number of shares held at time $t$ is
% \[
% \boxed{
% \;
% \Delta_t
% =
% \frac{\partial p}{\partial s}(t,S_t)
% =
% \mathbf N\!\left(
% d_+(t,S_t)
% \right),
% \;}
% \]
% where
% \[
% d_+(t,s)
% =
% \frac{\ln\!\left(\frac{s}{K e^{-r(T-t)}}\right)}{\sigma\sqrt{T-t}}
% +
% \frac12\sigma\sqrt{T-t}.
% \]
% \end{theorem}

% \begin{remark}[Economic meaning]
% \begin{itemize}[leftmargin=2em]
% \item $\Delta_t\in(0,1)$ represents the fraction of one share needed
% to hedge the option.
% \item When the option is deep in-the-money, $\Delta_t\approx 1$.
% \item When it is far out-of-the-money, $\Delta_t\approx 0$.
% \end{itemize}
% \end{remark}

% % ----------------------------------------------------------
% \subsection{Connection with PDE and martingale pricing}
% % ----------------------------------------------------------

% \begin{itemize}[leftmargin=2em]
% \item The Black--Scholes PDE comes from \emph{replication and delta hedging}.
% \item The expectation formula comes from \emph{martingale pricing under $\Q$}.
% \item The fact that both approaches lead to the same price is a
% manifestation of \textbf{no-arbitrage}.
% \end{itemize}

% \begin{remark}[Key takeaway]
% \[
% \text{Hedging (PDE)} \quad \Longleftrightarrow \quad
% \text{Pricing (expectation under }\Q\text{)}.
% \]
% The drift $\mu$ of the real-world dynamics never appears in the final price.
% \end{remark}

% % ==========================================================
% \section{Material required for the assignment}
% % ==========================================================

% The following results and tools are assumed and may be used freely.

% \begin{itemize}[leftmargin=2em]

% \item \textbf{Brownian motion.}
% Independent increments, Gaussianity, scaling, and conditional distributions:
% \[
% W_T-W_t \sim \mathcal N(0,T-t), \quad \text{independent of } \mathcal F_t.
% \]

% \item \textbf{It\^o's formula (1D).}
% For $X_t$ solving
% \(
% \dd X_t=\mu_t\,\dd t+\sigma_t\,\dd W_t
% \)
% and $f\in C^{1,2}$,
% \[
% \dd f(t,X_t)
% =
% \partial_t f\,\dd t
% +\partial_x f\,\dd X_t
% +\tfrac12\sigma_t^2\,\partial_{xx}f\,\dd t.
% \]

% \item \textbf{Geometric Brownian motion.}
% If
% \(
% \dd S_t=\mu S_t\,\dd t+\sigma S_t\,\dd W_t,
% \)
% then
% \[
% S_t
% =
% S_0\exp\!\left(
% (\mu-\tfrac12\sigma^2)t+\sigma W_t
% \right).
% \]

% \item \textbf{Black--Scholes PDE.}
% Derived from self-financing replication and delta-hedging:
% \[
% \partial_t p
% + rs\,\partial_s p
% + \tfrac12\sigma^2 s^2\,\partial_{ss}p
% - rp
% =0.
% \]

% \item \textbf{Heat equation.}
% Solution of
% \(
% \partial_\tau u=\tfrac12\sigma^2\partial_{yy}u
% \)
% with initial condition $u(0,y)=\phi(y)$:
% \[
% u(\tau,y)
% =
% \E\!\left[\phi(y+\sigma\sqrt{\tau}\,Z)\right].
% \]

% \item \textbf{Gaussian expectations.}
% If $Z\sim\mathcal N(0,1)$, then
% \[
% \E[e^{aZ}]=e^{\frac12 a^2},
% \qquad
% \E[\Phi(Z)]=\int \Phi(z)\,\varphi(z)\,\dd z.
% \]

% \item \textbf{Risk-neutral pricing (formal).}
% For European payoffs,
% \[
% V_t
% =
% \E^{\Q}\!\left[
% e^{-\int_t^T r\,\dd s}H
% \mid\mathcal F_t
% \right].
% \]

% \end{itemize}

% \begin{remark}
% All computations in the assignment rely only on these tools.
% No advanced measure-theoretic results beyond Girsanov's theorem
% are required.
% \end{remark}

% \subsection*{B.5 Greeks (definitions + interpretation)}
% Assume $p(0,S_0)$ is the time-0 option price.
% Define the standard sensitivities (Greeks):
% \[
% \Delta := \frac{\partial p}{\partial s},\qquad
% \Gamma := \frac{\partial^2 p}{\partial s^2},\qquad
% \Theta := \frac{\partial p}{\partial t},\qquad
% \text{Vega}:=\frac{\partial p}{\partial \sigma},\qquad
% \rho:=\frac{\partial p}{\partial r}.
% \]
% \begin{itemize}[leftmargin=1.5em, itemsep=0.2em]
% \item $\Delta$ = hedge ratio (number of shares in the replicating strategy).
% \item $\Gamma$ = curvature; drives rebalancing intensity (liquidity/transaction cost sensitivity).
% \item $\Theta$ = time decay.
% \item Vega = sensitivity to volatility misspecification.
% \item $\rho$ = sensitivity to the interest rate.
% \end{itemize}


% \begin{itemize}
% \item Connect to replication: $\Delta_t=\partial_s p(t,S_t)$ is \emph{the} hedge.
% \item Give intuition for $\Gamma>0$ for convex payoffs (calls/puts).
% \end{itemize}


% \subsection*{B.6 From PDE to an expectation (motivation)}
% We now explain why the PDE solution can be written as a discounted expectation
% under a measure where the stock drift is $r$.

% \paragraph{Heuristic step.}
% Consider an auxiliary process $\widetilde S$ defined by
% \[
% \frac{\dd \widetilde S_t}{\widetilde S_t}=r\,\dd t+\sigma\,\dd \widetilde W_t,
% \]
% for some Brownian motion $\widetilde W$ (under some probability measure $\Qbb$).
% Then It\^o + standard arguments (Feynman--Kac viewpoint) suggest that the PDE solution satisfies
% \[
% \boxed{\ p(t,s)=\E^{\Qbb}\!\Big[e^{-r(T-t)}\,g(\widetilde S_T)\,\Big|\,\widetilde S_t=s\Big].\ }
% \]
% In particular,
% \[
% p(0,S_0)=\E^{\Qbb}\!\big[e^{-rT}\,g(S_T)\big],
% \qquad \text{with}\quad \frac{\dd S_t}{S_t}=r\,\dd t+\sigma\,\dd \widetilde W_t.
% \]


% \begin{itemize}
% \item Tell students: ``We will justify this properly soon; today it's motivation.''
% \item Key message: pricing uses a measure where discounted $S$ has zero drift.
% \end{itemize}


% ==========================================================
% \section*{Part C --- Girsanov: make the drift become $r$}
% % ==========================================================

% \subsection*{C.1 Discounted stock and the drift we want to remove}
% Let $\bar S_t:=D_t S_t=e^{-rt}S_t$. Product rule (since $D$ has finite variation):
% \[
% \dd \bar S_t
% =
% \dd(e^{-rt}S_t)
% =
% e^{-rt}\dd S_t - r e^{-rt}S_t\,\dd t
% =
% (\mu-r)\bar S_t\,\dd t + \sigma \bar S_t\,\dd W_t.
% \]
% Define the market price of risk
% \[
% \theta:=\frac{\mu-r}{\sigma}.
% \]
% Then
% \[
% \dd \bar S_t = \sigma \bar S_t(\theta\,\dd t+\dd W_t).
% \]


% \begin{itemize}
% \item Compute $\dd(e^{-rt}S_t)$ explicitly.
% \item Introduce $\theta=(\mu-r)/\sigma$ and state: ``we want to eliminate $\theta\,dt$''.
% \end{itemize}


% \subsection*{C.2 Exponential martingale and change of measure}
% Define
% \[
% Z_t:=\exp\Big(-\theta W_t-\frac12\theta^2 t\Big).
% \]
% Assume (here automatic since $\theta$ constant) that $(Z_t)_{t\le T}$ is a true martingale, so $\E[Z_T]=1$.
% Define a new probability measure $\Qbb$ on $\F_T$ by
% \[
% \frac{\dd\Qbb}{\dd\Pbb}\Big|_{\F_T}=Z_T.
% \]
% Girsanov theorem implies that
% \[
% \widetilde W_t:=W_t+\theta t
% \]
% is a Brownian motion under $\Qbb$.


% \begin{itemize}
% \item Link to earlier example: exponential martingale ``kill the drift''.
% \item State Girsanov (no proof now): under $\Qbb$, $W$ shifts by $+\theta t$.
% \end{itemize}


% \subsection*{C.3 Apply to discounted stock (drift becomes zero)}
% Rewrite $\dd W_t = \dd\widetilde W_t-\theta\,\dd t$ and plug into
% $\dd\bar S_t=\sigma\bar S_t(\theta\,\dd t+\dd W_t)$:
% \[
% \dd\bar S_t=\sigma\bar S_t\big(\theta\,\dd t+\dd\widetilde W_t-\theta\,\dd t\big)
% =\sigma\bar S_t\,\dd\widetilde W_t.
% \]
% Hence, under $\Qbb$, $(\bar S_t)_{t\le T}$ is a martingale, i.e.
% \[
% \boxed{\ \dd(e^{-rt}S_t)=\sigma e^{-rt}S_t\,\dd\widetilde W_t.\ }
% \]
% Equivalently, under $\Qbb$:
% \[
% \boxed{\ \frac{\dd S_t}{S_t}=r\,\dd t+\sigma\,\dd\widetilde W_t.\ }
% \]


% \begin{itemize}
% \item The only change is the drift: $\mu\leadsto r$. Volatility unchanged.
% \item This is the rigorous bridge to the expectation formula.
% \end{itemize}

% ==========================================================
\section*{Part C --- Change of probability measure and Girsanov:
make the drift become $r$}
% ==========================================================

\subsection*{C.0 Why change the probability measure? (conceptual motivation)}

Up to now, pricing was obtained via \emph{hedging}, leading to the
Black--Scholes PDE.
We now explain why the same price can be written as a discounted expectation
under a suitable probability measure.

The key idea is the following:

\begin{quote}
\emph{We may reweight scenarios so that the discounted stock price has zero drift.}
\end{quote}

This reweighting is precisely what a \emph{change of probability measure} does.

\paragraph{Change of measure (basic definition).}
Let $(\Omega,\mathcal F,\Pbb)$ be a probability space.
Given a nonnegative random variable $Z$ such that $\E[Z]=1$, define a new probability
measure $\Qbb$ by
\[
\Qbb(A):=\E[\mathbf 1_A Z],\qquad A\in\mathcal F.
\]
We write
\[
Z=\frac{d\Qbb}{d\Pbb}
\quad\text{(Radon--Nikodym derivative).}
\]

\paragraph{Key identity.}
For any integrable random variable $X$,
\[
\E^{\Qbb}[X]=\E[XZ].
\]
If $(Z_t)$ is a positive martingale with $\E[Z_T]=1$ and
$\frac{d\Qbb}{d\Pbb}\big|_{\mathcal F_T}=Z_T$, then for $0\le s\le t\le T$,
\[
\E^{\Qbb}[Y\mid\mathcal F_s]
=
\frac{1}{Z_s}\,\E[YZ_t\mid\mathcal F_s].
\]


\textbf{On the board (intuition).}
\begin{itemize}
\item Changing measure = reweighting paths.
\item Large $Z(\omega)$ makes scenario $\omega$ more likely under $\Qbb$.
\item We will choose $Z$ so that the drift of the discounted stock disappears.
\end{itemize}


% ----------------------------------------------------------
\subsection*{C.1 Discounted stock and the drift we want to remove}
% ----------------------------------------------------------

Recall the Black--Scholes model under $\Pbb$:
\[
\frac{dS_t}{S_t}=\mu\,dt+\sigma\,dW_t,
\qquad
B_t=e^{rt},
\qquad
D_t=e^{-rt}.
\]
Define the discounted stock
\[
\bar S_t := D_t S_t = e^{-rt}S_t.
\]
Since $D$ has finite variation, the product rule gives
\[
d\bar S_t
=
e^{-rt}dS_t - r e^{-rt}S_t\,dt
=
(\mu-r)\bar S_t\,dt + \sigma \bar S_t\,dW_t.
\]

Introduce the \emph{market price of risk}
\[
\theta:=\frac{\mu-r}{\sigma}.
\]
Then
\[
d\bar S_t=\sigma \bar S_t(\theta\,dt+dW_t).
\]


\textbf{Key observation.}
\begin{itemize}
\item Under $\Pbb$, the discounted stock has drift $(\mu-r)$.
\item We want a measure under which this drift vanishes.
\end{itemize}


% ----------------------------------------------------------
\subsection*{C.2 Exponential martingale and change of measure}
% ----------------------------------------------------------

Define the exponential process
\[
Z_t:=\exp\Big(-\theta W_t-\tfrac12\theta^2 t\Big).
\]
Since $\theta$ is constant, $(Z_t)$ is a true martingale and $\E[Z_T]=1$.

Define a new probability measure $\Qbb$ on $\mathcal F_T$ by
\[
\frac{d\Qbb}{d\Pbb}\Big|_{\mathcal F_T}=Z_T.
\]

\begin{theorem}[Girsanov (constant drift case)]
Under $\Qbb$, the process
\[
\widetilde W_t:=W_t+\theta t
\]
is a standard Brownian motion.
\end{theorem}


\textbf{Interpretation.}
\begin{itemize}
\item Under $\Pbb$: $W$ is Brownian.
\item Under $\Qbb$: the Brownian motion is $\widetilde W=W+\theta t$.
\item Drift has moved from the SDE into the measure.
\end{itemize}


% ==========================================================
\subsection*{How to compute expectations under $\Qbb$ (examples)}
% ==========================================================

Recall the constant-drift Girsanov setup:
\[
Z_t=\exp\!\left(-\theta W_t-\frac12\theta^2 t\right),\qquad
\frac{d\Qbb}{d\Pbb}\Big|_{\F_T}=Z_T,
\qquad
\widetilde W_t:=W_t+\theta t \text{ is Brownian under }\Qbb.
\]

\begin{remark}[Change-of-measure identities]
For any integrable $\F_t$-measurable random variable $Y$,
\[
\E^\Qbb[Y]=\E^\Pbb[Y\,Z_t].
\]
More generally, for $0\le s\le t\le T$ and integrable $\F_t$-measurable $Y$,
\[
\E^\Qbb[Y\mid \F_s]=\frac{1}{Z_s}\,\E^\Pbb[Y\,Z_t\mid \F_s].
\]
\end{remark}

% ----------------------------------------------------------
\subsubsection*{Example 1: distribution of $W_t$ under $\Qbb$}
% ----------------------------------------------------------

Since $\widetilde W_t=W_t+\theta t$ is $\mathcal N(0,t)$ under $\Qbb$, we have
\[
W_t=\widetilde W_t-\theta t \sim \mathcal N(-\theta t,\,t)\quad \text{under }\Qbb.
\]
Hence for any bounded measurable $\varphi:\R\to\R$,
\[
\E^\Qbb[\varphi(W_t)]
=
\E\Big[\varphi(\widetilde W_t-\theta t)\Big]
=
\int_{\R}\varphi(x)\,
\frac{1}{\sqrt{2\pi t}}
\exp\!\left(-\frac{(x+\theta t)^2}{2t}\right)\,dx.
\]

In particular,
\[
\E^\Qbb[W_t]=-\theta t,\qquad \text{Var}^\Qbb(W_t)=t.
\]

% ----------------------------------------------------------
\subsubsection*{Example 2: exponential moments (sanity check)}
% ----------------------------------------------------------

Let $a\in\R$. Using the Gaussian law under $\Qbb$:
\[
\E^\Qbb\!\left[e^{aW_t}\right]
=
\exp\!\left(a(-\theta t)+\frac12 a^2 t\right)
=
\exp\!\left(\frac12 a^2 t-a\theta t\right).
\]

Equivalently (same computation via $\Pbb$),
\[
\E^\Qbb\!\left[e^{aW_t}\right]
=
\E^\Pbb\!\left[e^{aW_t}Z_t\right]
=
\E^\Pbb\!\left[\exp\!\left((a-\theta)W_t-\frac12\theta^2 t\right)\right]
=
\exp\!\left(\frac12(a-\theta)^2 t-\frac12\theta^2 t\right).
\]

% ----------------------------------------------------------
\subsubsection*{Example 3: geometric Brownian motion and $\E^\Qbb[f(S_t)]$}
% ----------------------------------------------------------

Consider a stock of the form
\[
S_t = S_0\exp\!\left(\alpha t+\sigma W_t\right),
\qquad \sigma>0.
\]
Under $\Qbb$, since $W_t=\widetilde W_t-\theta t$, we can rewrite
\[
S_t
=
S_0\exp\!\left((\alpha-\sigma\theta)t+\sigma \widetilde W_t\right).
\]
Hence, under $\Qbb$,
\[
\ln S_t \sim \mathcal N\!\left(\ln S_0+(\alpha-\sigma\theta)t,\ \sigma^2 t\right).
\]

Therefore for any measurable $f$ such that $f(S_t)$ is integrable,
\[
\boxed{
\;
\E^\Qbb[f(S_t)]
=
\E\Big[f\!\Big(S_0\exp\!\big((\alpha-\sigma\theta)t+\sigma\sqrt{t}\,Z\big)\Big)\Big],
\qquad Z\sim\mathcal N(0,1).
\;}
\]
Equivalently, in integral form,
\[
\E^\Qbb[f(S_t)]
=
\int_{\R} f\!\Big(S_0 e^{(\alpha-\sigma\theta)t+\sigma x}\Big)\,
\frac{1}{\sqrt{2\pi t}}e^{-x^2/(2t)}\,dx.
\]

% ----------------------------------------------------------
\subsubsection*{Example 4 (Black--Scholes drift switch): make the drift become $r$}
% ----------------------------------------------------------

Suppose under $\Pbb$,
\[
\frac{dS_t}{S_t}=\mu\,dt+\sigma\,dW_t,
\qquad \mu\in\R,\ \sigma>0.
\]
Fix a constant $r$ and choose
\[
\theta:=\frac{\mu-r}{\sigma}.
\]
Then under $\Qbb$ (since $dW_t=d\widetilde W_t-\theta\,dt$),
\[
\frac{dS_t}{S_t}
=
\mu\,dt+\sigma(d\widetilde W_t-\theta\,dt)
=
r\,dt+\sigma\,d\widetilde W_t.
\]
So
\[
S_t
=
S_0\exp\!\left(\Big(r-\frac12\sigma^2\Big)t+\sigma \widetilde W_t\right)
\quad \text{under }\Qbb.
\]

Thus for any integrable $f$,
\[
\boxed{
\;
\E^\Qbb[f(S_t)]
=
\E\Big[f\!\Big(S_0\exp\!\big((r-\tfrac12\sigma^2)t+\sigma\sqrt{t}\,Z\big)\Big)\Big],
\qquad Z\sim\mathcal N(0,1).
\;}
\]

% ----------------------------------------------------------
\subsubsection*{Example 5: conditional expectation (useful for pricing)}
% ----------------------------------------------------------

Under $\Qbb$, $\widetilde W_t-\widetilde W_s\sim\mathcal N(0,t-s)$ independent of $\F_s$.
Hence, for $s<t$,
\[
S_t
=
S_s\exp\!\left(\Big(r-\frac12\sigma^2\Big)(t-s)+\sigma(\widetilde W_t-\widetilde W_s)\right).
\]
Therefore, for any measurable $f$ with integrable $f(S_t)$,
\[
\boxed{
\;
\E^\Qbb[f(S_t)\mid \F_s]
=
\Psi_{t-s}(S_s),
\quad
\Psi_{u}(x):=\E\Big[f\!\big(x e^{(r-\frac12\sigma^2)u+\sigma\sqrt{u}\,Z}\big)\Big].
\;}
\]
This is the basic ``lognormal transition'' identity used in Black--Scholes computations.


% ----------------------------------------------------------
\subsection*{C.3 Apply to the discounted stock}
% ----------------------------------------------------------

Rewrite $dW_t=d\widetilde W_t-\theta\,dt$ and substitute into
\[
d\bar S_t=\sigma \bar S_t(\theta\,dt+dW_t).
\]
We obtain
\[
d\bar S_t=\sigma \bar S_t\,d\widetilde W_t.
\]

\[
\boxed{\quad (e^{-rt}S_t)_{t\le T}\ \text{is a martingale under }\Qbb.\quad}
\]

Equivalently, under $\Qbb$ the stock price satisfies
\[
\boxed{\quad \frac{dS_t}{S_t}=r\,dt+\sigma\,d\widetilde W_t.\quad}
\]


\textbf{Punchline.}
\begin{itemize}
\item Drift changes: $\mu \longrightarrow r$.
\item Volatility unchanged.
\item Discounted prices become martingales.
\end{itemize}


% ----------------------------------------------------------
\subsection*{C.4 Link with pricing}
% ----------------------------------------------------------

Since $e^{-rt}S_t$ is a martingale under $\Qbb$, it is natural to expect that
discounted derivative prices should also be martingales.
This leads to the representation
\[
p(t,s)
=
\E^{\Qbb}\!\left[e^{-r(T-t)}g(S_T)\mid S_t=s\right],
\]
which coincides with the solution of the Black--Scholes PDE.

\begin{remark}
Today this is presented as motivation.
A full justification via no-arbitrage and self-financing strategies
will be given in the next lecture.
\end{remark}


\end{document}